\begin{abstract}
\noindent
\textbf{Objective:}
To measure adherence to a national inpatient hyperglycemia guideline at the
level of its individual recommendations across a hospital network, to
determine whether it changed after the guideline was issued, and to identify
which recommendations track subsequent outcomes.

\noindent
\textbf{Materials and Methods:}
Retrospective multicenter study of 48,354 admissions among 29,577 adults
with diabetes or inpatient hyperglycemia at six Israeli general hospitals,
August 2022 to December 2025. Each of 14 recommendations was formalized as a
quality-assessment temporal pattern over knowledge-based temporal
abstractions of the record, scored wherever it applied. Pre--post adherence was compared by
Welch $t$-test with false discovery rate correction. Associations with 10
complications were estimated by bias-reduced logistic regression over 24-,
48-, and 72-hour windows, and by inverse probability of
treatment weighting contrasting above- with at-or-below-median adherence.

\noindent
\textbf{Results:}
Adherence ranged from near-universal admission and medication-management
actions to substantial gaps in insulin management (means, 0.15--0.28). Overall adherence rose modestly after issue (0.682 to 0.703,
$P<.001$), with marked inter-hospital heterogeneity. After weighting,
above-median adherence carried 44\% lower odds of 30-day mortality (odds
ratio [OR], 0.557; 95\% CI, 0.514--0.603), 82\% lower odds of hyperglycemia
(OR, 0.175) and 52\% lower odds of severe hyperglycemia (OR, 0.485), against
20\% higher odds of hypoglycemia (OR, 1.204).

\noindent
\textbf{Discussion:}
The least-implemented recommendations were largely those whose adherence
tracked outcomes: repeated insulin decisions spread across a stay, rather
than discrete actions on admission.

\noindent
\textbf{Conclusion:}
Guideline adherence can be assessed automatically and per-recommendation at
national scale, identifying where fuller implementation would most plausibly
matter.
\end{abstract}
